% ECONOMICS_PROBLEM: {"id": "C73-256", "title": "Stationary measures with common noise: historical question and remaining scope", "jel_code": "C73", "source_id": "OP-0209"}
% !TeX program = xelatex
\documentclass[11pt,a4paper]{article}
\usepackage[margin=23mm]{geometry}
\usepackage{fontspec}
\setmainfont{Latin Modern Roman}
\usepackage{amsmath,amssymb,mathtools}
\usepackage{xcolor,enumitem,booktabs,longtable,array,xurl,hyperref}
\definecolor{muted}{HTML}{53636C}
\hypersetup{colorlinks=true,urlcolor=blue,linkcolor=blue}
\setlength{\parindent}{0pt}
\setlength{\parskip}{5pt}
\newcommand{\R}{\mathbb R}
\newcommand{\E}{\mathbb E}
\newcommand{\N}{\mathbb N}
\newcommand{\ind}{\mathbf 1}
\newcommand{\Var}{\operatorname{Var}}
\newcommand{\Cov}{\operatorname{Cov}}
\newcommand{\supp}{\operatorname{supp}}
\DeclareMathOperator*{\argmin}{arg\,min}
\DeclareMathOperator*{\argmax}{arg\,max}
\newcommand{\entrymeta}[3]{{\sffamily\small\color{muted}\textbf{\texttt{#1}}\quad|\quad #2\par #3\par}}
\newcommand{\Problemref}[1]{\texttt{#1}}
\newcommand{\JELref}[2]{\texttt{#2} (JEL \texttt{#1})}
\begin{document}
\section*{Stationary measures with common noise: historical question and remaining scope}
\textbf{Problem ID:} \texttt{C73-256}\quad \textbf{JEL:} \texttt{C73}\par
% BEGIN ECONOMICS STATEMENT
\label{problem:OP-0209}
% GAME_THEORY_PROBLEM: {"local_id": "GT-079", "original_number": 79, "kind": "H", "entry_kind": "statement_only", "published": false, "review_required": true, "category": "game_theory"}
\label{game:GT-079}
{\sffamily\small\color{muted}
{\small\textbf{GT-079} | Mean-field games | H: the blanket historical status is superseded in a standard monotone setting\par}
\par}
\subsection*{Definitions and mathematical statement}
Fix $\beta>0$, a continuous running coupling $F:\mathbb T^d\times\mathcal P(\mathbb T^d)\to\mathbb R$, and independent $d$-dimensional Brownian motions $B,W$. A stationary Markov MFG equilibrium consists of a feedback $A:\mathbb T^d\times\mathcal P(\mathbb T^d)\to\mathbb R^d$, a law $\pi\in\mathcal P(\mathcal P(\mathbb T^d))$, and a stationary measure process $(\mu_t)$ such that
\[
 dX_t=A(X_t,\mu_t)dt+\sqrt2\,dB_t+\sqrt{2\beta}\,dW_t,
 \quad \mu_t=\mathcal L(X_t\mid\mathcal F_t^0),\quad
 \mathcal L(\mu_t)=\pi.
\]
Require $A$ to be bounded and Lipschitz in $(x,m)$, using the torus distance and the $1$-Wasserstein distance. The common filtration contains $\mu_0$ and $(W_s)_{s\leq t}$, with $\mu_0\sim\pi$ and $\mathcal L(X_0\mid\mu_0)=\mu_0$; $(\mu_0,X_0)$ is independent of future Brownian increments. Stationarity means invariance of every finite-dimensional law of $(\mu_t)$ under time translations, not that each realized $\mu_t$ is constant. The feedback must minimize, against the fixed common-noise environment $(\mu_t)$ and with the same initial-state law, the long-run cost
\[
 J(\alpha;\mu)=\limsup_{T\to\infty}\frac1T
 \mathbb E\!\int_0^T\left(\tfrac12|\alpha_t|^2+F(X_t^\alpha,\mu_t)\right)dt,
 \quad
 dX_t^\alpha=\alpha_tdt+\sqrt2\,dB_t+\sqrt{2\beta}\,dW_t.
\]
Controls are progressively measurable for the filtration generated by the initial variables and both noises, with $\mathbb E\int_0^T|\alpha_t|^2dt<\infty$ for every finite $T$. Since $F$ is continuous on the compact state/measure space, the limsup cost is defined, possibly $+\infty$. In weak measure notation the population equation is
\[
 d\mu_t=\big((1+\beta)\Delta\mu_t-\operatorname{div}(\mu_t A(\cdot,\mu_t))\big)dt
             -\sqrt{2\beta}\,D\mu_t\cdot dW_t.
\]
If this equation with fixed $A$ defines a Markov semigroup $P_t^A$, invariance means $\pi P_t^A=\pi$ for every $t\geq0$. Constructing an arbitrary drift's invariant law does not by itself establish the equilibrium optimality condition.

\textbf{Historical target.} Establish such a stationary regime for the smooth, Lasry--Lions-monotone, quadratic-cost, torus model with additive Brownian common noise. The general formulation above is editorial; the cited later theorem constructs a stationary MFG system and its ergodic master corrector in this standard setting.

\textbf{Remaining agenda.} For a precisely chosen nonmonotone coupling class $\mathcal C$, determine which $F\in\mathcal C$ admit a stationary equilibrium law, and when that law is unique. No universal existence or uniqueness assertion for arbitrary nonmonotone $F$ is made here.
\subsection*{Short English statement}
Does the random population distribution have an equilibrium invariant law, and under which assumptions is it unique?
\subsection*{Variants and scope boundaries}
The compact torus, additive common noise, and monotonicity matter. Noncompact states require confinement; nonmonotone games may require an equilibrium selection. A stationary law of $(Du,\mu)$ and one of the entire value $u$ have different normalizations.
\subsection*{Source relationship and status}
The 2022 chapter's blanket status is outdated: the 2025 preprint constructs a monotone stationary regime. Broader extensions remain an agenda, not a verified universal open conjecture.
\subsection*{References}
\begingroup\footnotesize\raggedright
{\footnotesize\raggedright
\textbf{[S1]} Cardaliaguet and Delarue. \emph{Selected topics in mean field games}, ICM2022, vol.~5, pp.~3660--3703; DOI: 10.4171/ICM2022/17. \textit{Locator:} Section~3.2, final paragraph, p.~3680.\par
\url{https://ems.press/content/book-chapter-files/33265}\par\smallskip
\textbf{[S2]} Pierre Cardaliaguet, Rapha\"el Maillet, and Wenbin Yan (2025). \emph{On the long-time behavior of mean field game systems with a common noise}. arXiv:2509.17443v1, 22 September 2025; preprint.\par
\textit{Locators:} assumptions (Hf), (Hg), (Hm), Section~2, pp.~6--8; Theorem~3.9, p.~14; Theorem~6.6 and (6.11), p.~44, stationary measure process. The preprint's proofs have not been independently audited here.\par
\url{https://arxiv.org/pdf/2509.17443v1}\par}
\par\endgroup

\subsection*{Common conventions: Source mathematical conventions}

\textbf{Finite games.} For a finite set $A$, $\Delta(A)$ is its probability
simplex. Players choose independent mixed strategies in Nash equilibrium;
correlation is allowed only when explicitly part of the solution concept.
In a game with payoff functions $u_i$ and strategy profile $x$, an additive
$\varepsilon$ Nash equilibrium satisfies
\[
 u_i(x)\geq u_i(x_i',x_{-i})-\varepsilon
 \quad\text{for every player $i$ and every admissible deviation $x_i'$}.
\]
For numerical approximation, payoffs are normalized as locally specified.
An $\varepsilon$ payoff residual is distinct from distance at most
$\varepsilon$ to the set of exact equilibria. Point convergence, convergence
of empirical play, and convergence in distance to an equilibrium set are
also distinct.

\textbf{Computational representations.} Unless stated otherwise, a complexity
question takes rational data with binary encoding; polynomial time is measured
in total bit length. A fixed-accuracy polynomial algorithm is not necessarily
polynomial in $1/\varepsilon$, and neither is necessarily polynomial in
$\log(1/\varepsilon)$. Succinct games and value, demand, or sampling oracles
must declare their interfaces. Exact real solutions require an explicit
representation; an unspecified list of arbitrary real numbers is not a
finite computational output. A claimed algorithmic barrier is meaningful
only for its specified model and reductions.

\textbf{Dynamic games.} Histories, information, observation, and allowable
randomization are part of the model. A stationary strategy depends only on
the current state; a positional strategy in a deterministic graph game
chooses one action at each controlled vertex. Nash equilibrium at an initial
state and equilibrium after every history are different requirements.
An expectation of a pathwise $\liminf$ payoff, a $\liminf$ of expected averages,
a limiting expected average, and a uniform finite-horizon equilibrium are
different criteria. None is silently substituted for another.

\textbf{Allocation and mechanisms.} A complete allocation assigns every item
exactly once unless otherwise stated. Zero-valued goods, negative values,
capacity constraints, feasibility, lotteries, and copies of goods affect
fairness definitions and are specified locally. Dominant-strategy incentive
compatibility, Bayesian incentive compatibility, universal truthfulness,
and truthfulness in expectation impose different inequalities. Whenever
an objective ratio has zero denominator, its convention is stated or those
instances are excluded.

\textbf{Infinite-dimensional models.} Expectations, integrals, stochastic
equations, and optimization objectives require their local regularity,
integrability, filtrations, and admissible domains. A backward--forward
mean-field system is a boundary-value problem; it is not an initial-value
semiflow without an additional construction. Long-time, large-population,
and vanishing-noise limits require an order of limits and a convergence
topology. If a minimum need not be attained, the objective is its infimum
or supremum and the approximation requirement is stated separately.
% END ECONOMICS STATEMENT
\end{document}
